A computer can perform {\em additions} $(15 + 3) = 18$, {\em subtractions}
$(15 - 3) = 12$,
{\em multiplications} $(15 \ast 3) = 45$ and {\em divisions} of integers $(15 \;/\!/\; 3) = 5$.
If the division does not come out exactly,
the result is the integer quotient. The remainder is given
by the {\em modulo} operation, written {\idt \%}. For example
$(17 \;/\!/\; 3) = 5$ and $(17 \;${\idt \%}$\; 3) = 2$ because dividing the dividend $17$
by the divisor $3$ gives the quotient $5$ and the remainder $2$. In mathematics a cross
`$\times$' is used to denote multiplication, whereas in \python\ the asterisk $\ast$ is used.
This avoids confusing the multiplication sign `$\times$' with the lowercase letter {\idt x}.
Likewise, two slashes `{\idt //}' are used for integer division rather than the sign
`$\div$', which is not found on computer keyboards.
Be careful not to write a single slash `{\idt /}': in \python, it does not give the
integer quotient ($17 \;/\; 3$ is not $5$). We will see later that it is used to compute with
decimal numbers.

To do a calculation involving several operations, we use pa\-ren\-the\-ses. For example
$((2 \ast 3) - 1)$ is equal to 5. The calculation begins by evaluating the innermost
expression in parentheses.
For example $((((5 + 1) \ast 3) - 4) \;/\!/\; 2) = (((6 \ast 3) - 4) \;/\!/\; 2) =
((18 - 4) \;/\!/\; 2) = (14 \;/\!/\; 2) = 7$.

\index{Integer expression}In these {\em integer expressions}
we can also use the value of the counters of ``for'' loops. For example,
during the successive iterations in the body of the loop
``{\mcl for}\ {\idt i}\ {\mcl in}\ {\idt range(1,\ 6):}\ \ldots'',
the counter {\idt i} is $1$, then $2$, then $3$, then $4$ and finally $5$.
In this loop body the expression $((2 \ast ${\idt  i}$) - 1)$ therefore successively has the value
$((2 \ast 1) - 1) = 1$, $((2 \ast 2) - 1) = 3$, $((2 \ast 3) - 1) = 5$, $((2 \ast 4) - 1) = 7$ and finally
$((2 \ast 5) - 1) = 9$.

\begin{exemple}[Pyramid]
To draw the pyramid below, we can start from the top and draw, from
top to bottom, successive rectangles of lengths 1, 3, 5, 7, 9. The length of the
side of the {\idt i}$^{\rm th}$ rectangle counting from the top is
therefore $((2\;\ast\;${\idt  i}$) - 1)$. In the program, this expression is written
`{\idt 2\ *\ i\ -\ 1}' between the parentheses of {\idt range} to give the number of iterations
of the loop that draws one side of the rectangle.

\pagebreak[3]
\vspace*{22mm}%
\includePostScript{Pyramide}{0.95}{0.95}%
\vspace*{-22mm}%
\input python/Pyramid.tex
\end{exemple}

%\pagebreak
\needspace{24\baselineskip}\noindent \rule{137mm}{0.1mm}\par\nopagebreak %---------------------------------------------
\begin{exercice}{ Complete the program {\idt SawtoothStakes}, replacing the
ellipsis `{\idt ...}' by an integer expression, so as to
ob\-tain the row of stakes shown below:

\raggedright{
\vspace*{20mm}%
\includePostScript{PiquetsEnDentsDeScie}{0.95}{0.95}\ \\%
}


\input python/SawtoothStakes.tex
\noindent $\Box$}\end{exercice}
\begin{exercice}{ Write \python\ programs commanding the \robot\ to produce the
geometric drawings of figure  \ref{fig:spcf} (page  \pageref{fig:spcf}). $\Box$}\end{exercice}
\begin{exercice}{{\bf (Difficult)} Complete the program {\idt WavesOfStakes} to
draw the row of stakes shown below:

\pagebreak[3]
\raggedright{
\vspace{0.5cm}

\includevspacePostScript[2.0cm]{VaguesDePiquets}{0.95}{0.95}
}

\input python/WavesOfStakes.tex
\noindent $\Box$}\end{exercice}
\begin{figure}[hb]%[htbp]
 \includevspacePostScript[9.3cm]{SpiralesEtQuad}{0.95}{0.95}
\caption{Greek key, spirals, triangular tiling and fringed square}
 \label{fig:spcf}
\end{figure}

\clearpage

\begin{exercice}{{\bf (Difficult)}
\begin{enumerate}
\item Write a \python\ program to draw the tiling of Nicomachus\@%
\footnote{\ Nicomachus of Gerasa, Greek mathematician, c.\ 1st century AD.}
shown in figure
\ref{fig:Nicomaque}  page \pageref{fig:Nicomaque} using dependent nested
``for'' loops. Describe the loop invariants
by comments in the program. The tiling of Nicomachus is made of 10
rings consisting of squares whose side lengths
are respectively 1, 2, 3, 4, \ldots, 10.

\item Compute the area of the square in two different ways
suggested by figure \ref{fig:Nicomaque2}
page \pageref{fig:Nicomaque2}:
 \begin{enumerate}
  \item by computing the length of the half-side of the square, obtained by adding
        those of the squares of the tiling,
  \item by computing the area of the quarter of the large square delimited by
   the diagonals of this large square (see this part in green in
   figure \ref{fig:Nicomaque2}, page \pageref{fig:Nicomaque2}), the area of this quarter of the large square
   being obtained by adding the areas of the small squares that make it up.
 \end{enumerate}
\item Deduce the equality:
\[ 1^{3} + 2^{3} + 3^{3} + 4^{3} + \ldots + 10^{3} = (1 + 2 + 3 + 4 + \ldots + 10)^{2}\]
\end{enumerate}
\noindent $\Box$}\end{exercice}
\begin{figure}[htbp]
 \includevspacePostScript[11.7cm]{Nicomaque}{0.95}{0.95}
 \caption{Tiling of Nicomachus (Scale $\frac{1}{4})$}
 \label{fig:Nicomaque}
\end{figure}

\begin{corrige} {The missing integer expression in the program {\idt Saw\-tooth\-Stakes} is:
\begin{center}
$(1 + ((${\idt i}$ - 1)\;${\idt \%}$\;5))$
\end{center}
which successively takes the values $1$, $2$, $3$, $4$, $5$, $1$, $2$, $3$, \dots when {\idt i} takes the
values $1$, $2$, $3$, $4$, $5$, $6$, $7$, $8$, \ldots
\noindent $\Box$}\end{corrige}
\begin{corrige} {We give one program; the others are in the folder
{\idt book/\allowbreak 08\_integer\_expressions}:
\input python/TrapezoidalSpiral.tex
\noindent $\Box$}\end{corrige}
\begin{corrige} {The missing integer expression in the program {\idt WavesOfStakes} is the following:
\begin{center}
$(((6-${\idt j}$)\ast ((${\idt i}$+1)\;${\idt \%}$\;2))+(${\idt j}$\ast(${\idt i}$\;${\idt \%}$\;2)))$
\end{center}
When {\idt i} is odd, $((${\idt i}$+1)\;${\idt \%}$\;2) = 0$ and $(${\idt i}$\;${\idt \%}$\;2) = 1$,
so the expression is equal to {\idt j}. When {\idt i} is even, $((${\idt i}$+1)\;${\idt \%}$\;2) = 1$ and
$(${\idt i}$\;${\idt \%}$\;2) = 0$, so the expression is equal to $(6-${\idt j}$)$.
$\Box$}\end{corrige}
\begin{corrige}{\hfill
\input python/NicomachusTiling.tex
\begin{figure}[htbp]
 \includevspacePostScript[12.1cm]{Nicomaque2}{0.95}{0.95}
 \caption{Side and area of the square of Nicomachus (Scale $\frac{1}{4}$)}
 \label{fig:Nicomaque2}
\end{figure}
If we divide one side of the large square in two, the length of this
half-side is the sum $(1 + 2 + \ldots + 10)$ of the lengths of the squares
of sizes 1, 2, 3, \ldots, 10 of the tiling. This can be seen clearly in
figure \ref{fig:Nicomaque2}  page \pageref{fig:Nicomaque2} thanks to the
ruler that we have placed below and to the right of the large square.
The area of the large square is therefore the square of the length of its side, that is:
\[(2 \times (1 + 2 + \ldots + 10))^{2} = 4 \times (1 + 2 + \ldots + 10)^{2}\]

If we divide the large square along its diagonals, we obtain a
triangle, as shown in green in
figure \ref{fig:Nicomaque2} (page \pageref{fig:Nicomaque2}). This triangle
contains 1 square of side length 1, 2 squares of side
length 2, \ldots, 10 squares of side length 10. Its area is
therefore $1 \times 1^{2} + 2 \times 2^{2} + \ldots + 10  \times  10^{2}$, that is,
$1^{3} + 2^{3} + \ldots + 10^{3}$. By turning this triangle 4 times, we
cover the large square completely. The area of the large square is
therefore:
\[4 \times (1^{3} + 2^{3} + \ldots + 10^{3})\]

The result is the same whichever way we compute the area
of the square of Nicomachus. Consequently:
\[4 \times (1^{3} + 2^{3} + \ldots + 10^{3}) = 4 \times (1 + 2 + \ldots + 10)^{2}\]
that is:
\[ 1^{3} + 2^{3} + \ldots + 10^{3} = (1 + 2 + \ldots + 10)^{2}\]

Another proof of this identity was given by Al-Karaji\/\@\footnote{\ Al-Karaji,
late 10th -- early 11th century, Persian mathematician.} and relies on tiling the square with
L-shaped pieces (gnomons), as illustrated below:

\begin{figure}[htbp]
 \includevspacePostScript[6.3cm]{AlKaragi}{0.95}{0.95}
 \caption{Tiling of Al-Karaji}
 %\label{fig:AlKaragi}
\end{figure}

\noindent The area of the outer L-shape with short side 10 is:
\begin{center}
$\begin{array}{ll}
(2 \times 10 \times (1 + 2 + \ldots + 10))  - 10^{2} & = 2 \times 10 \times \frac{10 \times (10 +1)}{2} - 10^{2}\\
& = -10^{3} + 10^{2} - 10^{2} = 10^{3}
\end{array}$
\end{center}
The area of the L-shape with short side 9 is $9^{3}$, that of the L-shape with short side 8 is $8^{3}$,  \ldots,
the area of the small square is $1 = 1^{3}$.
Consequently, the area of the large square, which
is the sum of the areas of the small square and of the L-shapes, is $1^{3} + 2^{3} + \ldots + 10^{3}$.
On the other hand, the length of the side of the large square is $(1 + 2 + \ldots + 10)$, so its area is
$(1 + 2 + \ldots + 10)^{2}$. Since the area of the large square is the same whichever way
we compute it, we obtain:
\[ 1^{3} + 2^{3} + \ldots + 10^{3} = (1 + 2 + \ldots + 10)^{2}\]
To justify that $(1 + 2 + \ldots + 10) = \frac{10 \times (10 +1)}{2}$, we use a rectangle of
width 10 and length (10 +1), half of which is tiled as follows:

\begin{figure}[htbp]
 \includevspacePostScript[4.3cm]{S1p2p3Etc}{0.95}{0.95}
 \caption{Computing $1 + 2 + 3 + \ldots + 10$}
 %\label{fig:1+2+3etc}
\end{figure}

Another geometric justification of the equality
$1^{3} + 2^{3} + \ldots + 10^{3} = (1 + 2 + \ldots + 10)^{2}$, proposed by S. Go\-lomb
\cite{golomb1}, is illustrated below:
%à la figure \ref{fig:golomb} (page \pageref{fig:golomb}).

\begin{figure}[htbp]
 \includevspacePostScript[6.0cm]{Golomb}{0.95}{0.95}
 \caption{Tiling of Golomb}
 \label{fig:golomb}
\end{figure}

\noindent The length of the side of the large square is indeed $(1 + 2 + \ldots + 10)$, so its area is
$(1 + 2 + \ldots + 10)^{2}$. This area is equal to the sum of the areas of the squares
making up the tiling, that is, 1 square of side 1, 2 squares of side 2, \ldots, 10
squares of side 10. Note
that for the squares with an even side length, shown in yellow in the figure,
there is an
overlap of two small squares, shown in red. But each of these small squares
is adjacent to an uncovered region of the same size, shown in white. We can therefore
remove one of the two overlapping small squares to fill the hole and obtain a perfect tiling.
 The area is therefore equal to $1^{3} + 2^{3} + \ldots + 10^{3}$.
\noindent $\Box$}\end{corrige}

%%%%%%%%%%%%%%%%%%%%%%%%%%%%%%%%%%%%%%%%%%%%%%%%%%%%%%%%%%%%%%%%%%%%%%%%%%%
\noindent \framebox[\textwidth]{
\begin{minipage}{\myparboxwidth}
\setlength{\parindent}{2em}
\small \sf
\vspace{\abovedisplayskip}
\section*{Addenda}
\indent \index{Integer expression}To make learning the syntax of \python\ easier, we write integer
expressions with parentheses.  {\em Unary} integer expressions (with
a single {\em operand} $E$) have the form $(-E)$ where $-$ is the minus sign {\em operator}.
{\em Binary} integer expressions (with a left {\em operand} $E_1$ and a right
operand $E_2$) have the form $(E_1\,$ {\mcl op} $\,E_2)$ where {\mcl op} is one of the
{\em operators} of addition $+$, subtraction $-$, multiplication $\ast$, integer
division {\idt //} or remainder of integer division {\idt \%} (modulo).
The operands $E$, $E_1$ and $E_2$ are an integer con\-stant (positive such as $1$, $17345$
(without blanks), \ldots
or negative such as $-10$, $-5247$, \ldots), a variable (for now a ``for'' loop
counter such as {\idt i}) or an integer
{\em subexpression} of the form we have just described.

Computing the value of an expression is called {\em evaluating} the expression. The rules
for evaluating fully parenthesized expressions are very simple. The value of a
constant is the integer denoted by this constant. The value of a variable {\idt i} used
as a counter in the body $C$ of a loop {\mcl for}\ {\idt i}\ {\mcl in}\ {\idt range(1,\ }$n${\idt \ +\ 1):}\ $C$
(where $n \geq 1$ is an integer constant) is successively 1, then 2,
then \dots, then $n$ during the successive iterations of the loop. The value of an expression $(-E)$
is obtained by computing the value $v$ of the expression $E$ and then applying the minus sign to it. The
value of an expression $(E_1\,$ {\mcl op} $\,E_2)$ is obtained by first computing the value
$v_1$ of the expression $E_1$, then computing the value $v_2$ of the expression $E_2$, then
computing the value $v_1\,$ {\mcl op} $\,v_2$ of the result. The same evaluation rules
apply to the subexpressions $E$, $E_1$ and $E_2$.

\indent To avoid cumbersome notation, in \python\ we can follow the common practice in
mathematics of using {\em priority rules} (also called {\em precedence
rules}) that allow us to remove
some redundant parentheses.
The evaluation rules for integer expressions are
defined and illustrated below:
\begin{list}{$\scriptscriptstyle\bullet$}{
}
\item In \python, as in mathematics, the integer expression
$10 - 10 - 10$ is equal to $((10 - 10) - 10)$, that is, $-10$.  Likewise $ 2 \ast 7 \,/\!/
\, 3\,${\idt \%}$\, 2$ is equal to $(((2 \ast 7) \,/\!/\, 3)\,${\idt \%}$\, 2)$, that is, $0$.
In both cases the rule is that operators of the same {\em priority} (namely $+$ and
$-$ on the one hand; $\ast$, {\idt //} and {\idt \%} on the other) are evaluated from left to right.
\item The expression $8 + - 2 \ast 3$ is equal to $(8 + ((-2) \ast 3))$, that is, $2$.  Likewise
\hbox{$1 - 5\,/\!/\,2$}
is equal to $(1 - (5\,/\!/\,2))$, that is, $-1$. The rule is that when there are operators of
different priorities, the operators of highest priority are evaluated first.
\end{list}

\vspace{\abovedisplayskip}
\end{minipage}
}

%%%%%%%%%%%%%%%%%%%%%%%%%%%%%%%%%%%%%%%%%%%%%%%%%%%%%%%%%%%%%%%%%%%%%%%%%%%%
%%%%%%%%%%%%%%%%%%%%%%%%%%%%%%%%%%%%%%%%%%%%%%%%%%%%%%%%%%%%%%%%%%%%%%%%%%%
\noindent \framebox[\textwidth]{
\begin{minipage}{\myparboxwidth}
\setlength{\parindent}{2em}
\small \sf
\vspace{\abovedisplayskip}
\noindent In \python\ the priorities are as follows:

\vspace{-6pt}

\begin{center}
\hspace{1.0cm}\begin{tabular}{|l|l|l|} \hline
\multicolumn{3}{|l|}{\em Priority of integer operators} \\ \hline \hline
$-$ (unary)                                & unary operators & high priority \\ \hline
$\ast$\ {\idt //}\ {\idt \%} & multiplicative operators & medium priority \\ \hline
$+$, $-$ (binary) & additive operators & low priority \\ \hline
\end{tabular}
\end{center}

\vspace{-6pt}

\begin{list}{$\scriptscriptstyle\bullet$}{
}
\item The expression $(8 + - 2) \ast 3$ is equal to $((8 + (- 2)) \ast 3)$, that is, $18$.  Likewise
\hbox{$(1 - 5)\,/\!/\,2$} is equal to $((1 - 5)\,/\!/\,2)$, that is, $-2$. In both cases we
begin by evaluating the subexpressions in parentheses.
\item When the division of a negative number does not come out exactly, \python\ rounds the
quotient down, that is, towards the smaller integer: $-7\,/\!/\,2$ is equal to $-4$
(and not $-3$ as in \ocaml\ or in \pascal) and the remainder $-7\;${\idt \%}$\;2$ is equal to $1$
(and not $-1$ as in \ocaml), so that we always have
$((-7\,/\!/\,2) \ast 2) + (-7\;${\idt \%}$\;2) = -7$. With a positive divisor, the remainder
{\idt \%} is therefore always positive or zero.
\end{list}

\smallskip

\noindent Equivalently, the order of evaluation of expressions can be defined
by stating to which operators the operands of an expression are bound:
\begin{list}{$\scriptscriptstyle\bullet$}{
}
\item An operand appearing between two operators of different priorities is bound to the operator
of higher priority (so $2 + 3 \ast 5 = (2 + (3 \ast 5))$).
\item An operand appearing between two operators of the same priority is bound to the operator
on the left (so $2 - 3 - 5 = ((2 - 3) - 5)$).
\item Expressions in parentheses are evaluated before being used as operands.
\end{list}

\smallskip

The rules for evaluating expressions change from one programming language to another.
For example, in {\sc apl} \cite{iverson1}, there are no operator priorities (because there are
too many operators in the language to hope to remember them) and expressions are right-associative,
so that $10 - 10 - 10 = (10 - (10 - 10)) = (10 - 0) = 10$. This rule, contrary
to mathematical tradition, has the advantage that the grammatical analysis of expressions is greatly
simplified for the computer.

Programming languages more classical than {\sc apl} fol\-low the mathematical
convention of left-to-right evaluation with priorities.


However, some languages have a great many operators.
For example, the {\sc C} language \cite{kernighan1} has 43 different operators divided into
15 lev\-els of pri\-or\-i\-ty.
The priorities of the operators are therefore not easy to remember.

\vspace{\abovedisplayskip}
\end{minipage}
}

%%%%%%%%%%%%%%%%%%%%%%%%%%%%%%%%%%%%%%%%%%%%%%%%%%%%%%%%%%%%%%%%%%%%%%%%%%%%
%%%%%%%%%%%%%%%%%%%%%%%%%%%%%%%%%%%%%%%%%%%%%%%%%%%%%%%%%%%%%%%%%%%%%%%%%%%
\noindent \framebox[\textwidth]{
\begin{minipage}{\myparboxwidth}
\setlength{\parindent}{2em}
\small \sf
\vspace{\abovedisplayskip}
Moreover, the rules change from one language to
another, which makes life difficult for computer scientists
who use several languages every day. To avoid mistakes they use redundant
parentheses,
which is what we advise beginners to do.

Finally, note that in \python\ integers are not limited: we can compute
with integers of several hundred digits, the only limit being the size of the
computer's memory. This is not the case in all languages. In \pascal\ on the first
personal computers, integers had to lie between $-32768$ and $32767$. In \ocaml,
they lie between {\idt min\_int} $= -4611686018427387904$ and
{\idt max\_int} $= 4611686018427387903$ and beyond that, the calculation gives a
wrong result without warning: $(${\idt max\_int}$ + 1)$ is equal to {\idt min\_int}!

\vspace{\abovedisplayskip}
\end{minipage}
}
%%%%%%%%%%%%%%%%%%%%%%%%%%%%%%%%%%%%%%%%%%%%%%%%%%%%%%%%%%%%%%%%%%%%%%%%%%%%
